% ── 7. Well-Architected Framework Alignment ───────────────────────────────────
\section{Well-Architected Framework Alignment}

The AWS Well-Architected Framework \cite{aws-well-architected} provides a structured lens for evaluating cloud architectures across six pillars: Operational Excellence, Security, Reliability, Performance Efficiency, Cost Optimization, and Sustainability. The sections below assess the LightsPlease platform against each pillar, identifying where the design aligns with best practices, where deliberate tradeoffs were made, and where gaps exist that should be addressed as the platform matures.

\subsection{Operational Excellence}

Operational excellence centers on the ability to run and monitor systems to deliver business value and to continually improve supporting processes and procedures. The architecture supports operational excellence primarily through its reliance on managed services, which offload the operational burden of patching, scaling, and availability monitoring for the underlying infrastructure. Amazon CloudWatch provides unified observability across all platform services, and alarms are connected to SNS topics for operational notification. AWS IoT Device Management enables over-the-air firmware updates to the device fleet, which is essential for long-term operational health of a distributed hardware product.

The platform should implement structured runbooks for common operational events — device connectivity degradation, elevated Lambda error rates, KVS ingest failures — and store these alongside the codebase so that on-call engineers have clear remediation paths. Infrastructure-as-code deployment reduces the operational risk of configuration drift between environments. Feature flags or canary deployments via Lambda weighted aliases allow new function versions to be validated under real traffic before full promotion, reducing the blast radius of defective releases.

% TODO: Expand with CI/CD pipeline specifics and observability dashboard design once those are defined.

\subsection{Security}

The Security pillar covers the protection of data, systems, and assets through identity and access management, threat detection, infrastructure protection, data security, and incident response. The architecture addresses security at each layer. Identity for consumers is managed through Amazon Cognito, which enforces password policies, supports multi-factor authentication, and issues short-lived tokens validated at the API Gateway boundary. Device identity is managed through X.509 certificates issued by AWS IoT Core's built-in certificate authority; each device has a unique credential and communicates over TLS. AWS IoT Core policies restrict what each device may publish to and subscribe from, preventing lateral movement between household namespaces.

AWS WAF protects the public-facing CloudFront distribution and API Gateway endpoint from common web exploits and volumetric attack patterns. All data at rest — in DynamoDB, S3, Timestream, and CloudWatch Logs — is encrypted using AWS-managed keys. IAM roles follow the principle of least privilege: each Lambda function is assigned a role with permissions scoped to only the specific resources it needs to access. Cross-account access is not required at this stage of the platform's maturity.

A key residual security risk is the potential for a device compromise to be used as a foothold within the IoT namespace. The mitigation is to scope each device's IoT policy to the minimum topic paths required for its function, so that a compromised device cannot publish to topics belonging to other devices or subscribe to control channels outside its own namespace. Certificate revocation procedures should be documented and tested as part of the incident response plan.

\subsection{Reliability}

Reliability addresses the ability of a workload to perform its intended function correctly and consistently. The architecture achieves reliability by building on services with strong durability and availability guarantees and by designing for asynchronous, decoupled flows that tolerate transient failures. DynamoDB provides 99.999\% availability for global tables and replicates data synchronously across multiple availability zones. S3 provides eleven nines of object durability. Lambda retries asynchronously invoked functions automatically on failure, and Step Functions provides durable orchestration with configurable retry and catch semantics for multi-step workflows.

The IoT message path is designed to tolerate intermittent device connectivity. MQTT's quality-of-service levels and IoT Core's persistent session support ensure that messages published while a device is temporarily disconnected are buffered and delivered upon reconnection, up to the platform's configured retention window. The user-facing API path is stateless at the Lambda layer, meaning that any invocation failure results in a clean retry rather than a partially committed state change.

% TODO: Add discussion of disaster recovery objectives (RPO/RTO) once these are defined by the business.

\subsection{Performance Efficiency}

Performance Efficiency is about using computing resources efficiently to meet system requirements and maintaining that efficiency as demand changes. The architecture uses purpose-built services for each workload type: DynamoDB for low-latency key-value access, Timestream for time-series ingest and query, KVS for video streaming, and Lambda for event-driven compute. This alignment between workload characteristics and service capabilities avoids the performance penalties that arise from forcing a workload into an ill-suited service.

CloudFront's global edge network reduces latency for consumers accessing the web application from geographically diverse locations by serving cached static assets from the edge rather than from the origin region. API Gateway's response caching can be applied to endpoints that return slowly changing data — such as household configuration or device capability metadata — to reduce Lambda invocation overhead for repeated identical requests. Lambda provisioned concurrency can be applied to latency-sensitive functions to eliminate cold-start delay for the first invocation after a period of inactivity, at the cost of a higher per-hour charge.

\subsection{Cost Optimization}

Cost optimization is discussed in detail in Section~\ref{sec:cost}. From a Well-Architected perspective, the design follows the key cost optimization principles by using consumption-based pricing (Lambda on-demand, DynamoDB on-demand), avoiding idle resource costs, and matching the resource type to the workload. The primary area where the architecture accepts higher cost in exchange for capability is Kinesis Video Streams, which carries a higher storage cost than a direct-to-S3 archival model but enables the live streaming capability that is central to the product's camera use case.

As the platform scales, the cost optimization path involves migrating DynamoDB tables with predictable traffic to provisioned capacity with auto-scaling, evaluating S3 Intelligent-Tiering for camera archival objects, and monitoring Lambda function memory allocation against actual utilization to right-size function configurations.

\subsection{Sustainability}

The Sustainability pillar encourages minimizing the environmental impact of cloud workloads. The serverless and managed-service architecture inherently supports sustainability goals by maximizing resource utilization: Lambda execution environments are shared across many customers, and AWS manages bin-packing at the infrastructure level in a way that individual operators cannot replicate in dedicated infrastructure. The platform does not run idle EC2 instances or reserved database servers; compute is consumed only when requests are being processed.

The most significant sustainability consideration in this architecture is the continuous video ingestion path, which consumes bandwidth and storage continuously for camera-equipped households. Enabling motion-triggered recording by default rather than continuous recording — and making continuous recording an opt-in setting — would meaningfully reduce the data volume processed and stored by the platform, with corresponding reductions in both cost and energy consumption. This is a product decision rather than a purely architectural one, but the architecture is designed to support it through IoT event-triggered recording policies.

% TODO: Reference AWS's carbon footprint tool and sustainability commitments if needed for the report.
